%! Principal Component Analysis — Unit 7, Lesson 7.3 — Exit Ticket (Answer Key)
\documentclass[10pt]{article}
\usepackage{linalg-article}
\usepackage{linalg-key}   % pulls in -boxes; do NOT also load -boxes
\newcommand{\TallMath}[1]{$\displaystyle #1\rule[-1.4em]{0pt}{3.2em}$}
\newcommand{\uu}{\mathbf{u}}
\newcommand{\vv}{\mathbf{v}}
\newcommand{\zero}{\mathbf{0}}
\newcommand{\T}{^{\top}}

\begin{document}

\pageheader{Unit 7, Lesson 7.3}{Exit Ticket --- Answer Key}
\namedateperiod

\vspace{0.15in}

No notes. Show your reasoning. A data set of four points has already been \textbf{centered}:
$(2,3),(3,2),(-2,-3),(-3,-2)$. Stacked as the rows of $A$, they give
$A\T A=\begin{bmatrix} 26 & 24\\ 24 & 26 \end{bmatrix}$.

\begin{enumerate}[leftmargin=1.8em, itemsep=15pt]
  \item \textbf{Find the principal components.} Solve $\det(A\T A-\lambda I)=0$ for $\lambda_1,\lambda_2$, then find PC1 (the unit eigenvector for the larger $\lambda$).
        \ansline{$(26-\lambda)^2-576=\lambda^2-52\lambda+100=0\Rightarrow\lambda=50,2$. For $\lambda=50$: $\begin{bsmallmatrix} -24 & 24\\ 24 & -24 \end{bsmallmatrix}\vv=\zero\Rightarrow \text{PC1}=\vv_1=\tfrac1{\sqrt2}\begin{bsmallmatrix} 1\\1 \end{bsmallmatrix}$.}
  \item \textbf{How much spread on PC1?} The spread along a component is $\sigma^2=\lambda$. What fraction of the total spread does PC1 capture? Is summarizing each point by its PC1 coordinate a good idea here?
        \ansline{$\dfrac{50}{50+2}=\dfrac{50}{52}\approx96\%$. Yes --- PC1 holds almost all the spread, so one number per point (its PC1 coordinate) keeps about $96\%$ of the information.}
  \item \textbf{Explain.} (a) Why do we \emph{center} the data before finding principal components? (b) What does keeping only PC1 let us do with the data? \ansline{(a) Centering puts the cloud at the origin so the eigenvectors of $A\T A$ describe how the data \emph{spreads}, not where its mean sits. (b) It reduces the dimension --- each two-feature point becomes a single PC1 coordinate, a compact summary that keeps most of the spread.}
\end{enumerate}

\begin{teachernote}
Sort into ``got it / arithmetic slip / concept slip.'' Q1 is the Lesson~7.1 recipe used as PCA: $\lambda=50,2$
and PC1 $=\tfrac1{\sqrt2}(1,1)$ (watch for taking the eigenvector of the \emph{smaller} $\lambda$). Q2 is the
spread fraction $50/52\approx96\%$ --- students should link a high percentage to ``good summary.'' Q3 is the
target justification: center so directions mean spread; keeping PC1 is dimension reduction (two numbers become
one). If many miss Q3, open 7.4 with a 30-second recap --- centering and ``keep the biggest $\sigma$'' are the
ideas the whole SVD story rests on.
\end{teachernote}

\end{document}
